\documentclass[12pt,a4paper]{ctexart}
\usepackage[margin=1in]{geometry}
\usepackage{booktabs}
\usepackage{amsmath,amssymb}
\usepackage{enumitem}
\usepackage{graphicx}
\usepackage{xurl}
\usepackage[colorlinks=true,linkcolor=blue,urlcolor=blue]{hyperref}

\title{实验部分（中文版）}
\author{}
\date{\today}

\begin{document}
\maketitle

\section{实验}
\label{sec:experiments_cn}

\subsection{Overview}

本节围绕三个问题展开：1）在当前数据协议下，TCDP 是否显著优于共享线性基线；2）TCDP 的关键增益来自哪些组件；3）这些结论在不同容量配置下是否稳定。

\subsection{实验协议}

所有结果统一使用以下目标域划分协议：
\begin{itemize}[leftmargin=1.2em]
  \item 划分方式：\texttt{table\_group}；
  \item 比例：train:val$=1:5$；
  \item 当前阶段不启用独立测试集。
\end{itemize}
对应数据文件为：
\begin{center}
\texttt{paper\_materials/data/dataset\_table\_group\_train\_val\_no\_test.pkl}
\end{center}
因此，本节定位为模型选择与机理验证，不对最终跨分布泛化作结论性表述。

\subsection{实现细节}

除被消融项外，训练设置固定如下：
\begin{itemize}[leftmargin=1.2em]
  \item 训练脚本：\texttt{src/train\_balanced\_sampling\_sep\_mlp\_shared\_calib\_step5a\_feat\_opt.py}；
  \item 结构编码：冻结 HGAT 并缓存结构嵌入（\texttt{freeze\_hgat}）；
  \item 优化器：AdamW；学习率调度：plateau；
  \item \texttt{batch\_size=2048}，\texttt{learning\_rate=2e-3}，\texttt{weight\_decay=5e-5}；
  \item \texttt{mlp\_dropout=0.2}，\texttt{z\_noise\_std=0.01}；
  \item 早停：\texttt{patience=12}，\texttt{min\_delta=8e-4}；
  \item 验证频率：\texttt{val\_eval\_interval=5}。
\end{itemize}
所有主表均基于 5 个随机种子：
\[
\{9294,9295,9296,9297,9298\}.
\]

\subsection{评价指标与统计口径}

对每个 seed，读取日志中 \texttt{[Best]} 对应检查点（按验证损失选择），记录 \texttt{Train R\textsuperscript{2}@Best}、\texttt{Val R\textsuperscript{2}}、\texttt{Val Loss} 与 \texttt{Best Epoch}。回归指标定义为：
\begin{equation}
R^2 = 1-\frac{\sum_i (y_i-\hat{y}_i)^2}{\sum_i (y_i-\bar{y})^2},
\end{equation}
\begin{equation}
\mathrm{MSE} = \frac{1}{N}\sum_i (y_i-\hat{y}_i)^2.
\end{equation}
跨 $K=5$ 个 seed 的聚合采用：
\begin{equation}
\mu(m)=\frac{1}{K}\sum_{k=1}^K m_k,\qquad
\sigma(m)=\sqrt{\frac{1}{K-1}\sum_{k=1}^K (m_k-\mu(m))^2}.
\end{equation}
本文表格统一报告 mean$\pm$std。

\subsection{核心模型对比}

表~\ref{tab:core_compare_cn} 给出核心架构对比结果。

\begin{table}[htbp]
\centering
\small
\caption{核心模型对比（\texttt{table\_group}，train:val$=1:5$，无测试集）。}
\label{tab:core_compare_cn}
\begin{tabular}{lcccc}
\toprule
方法 & NSeed & Train $R^2$ & Val $R^2$ & Val Loss \\
\midrule
Baseline(shared linear) & 5 & $0.9024\pm0.0086$ & $0.8726\pm0.0067$ & $0.1216\pm0.0064$ \\
TCDP(16,96) & 5 & $0.9586\pm0.0085$ & $0.9065\pm0.0168$ & $0.0892\pm0.0161$ \\
TCDP(32,64) & 5 & $0.9582\pm0.0216$ & $\mathbf{0.9234\pm0.0175}$ & $\mathbf{0.0731\pm0.0168}$ \\
\bottomrule
\end{tabular}
\end{table}

相较基线，TCDP(32,64) 的 Val $R^2$ 绝对提升 +0.0508，Val Loss 下降约 39.9\%。相较 TCDP(16,96)，TCDP(32,64) 的 Val $R^2$ 进一步提升 +0.0169，Val Loss 下降约 18.1\%。

\subsection{One-factor 消融（固定 TCDP(32,64)）}

采用 one-factor 设计：每次仅改动一个组件，其余配置固定。结果见表~\ref{tab:t32_ablation_cn}。

\begin{table}[htbp]
\centering
\scriptsize
\setlength{\tabcolsep}{3.5pt}
\caption{TCDP(32,64) 的 one-factor 消融（每项 5 seeds）。}
\label{tab:t32_ablation_cn}
\resizebox{\columnwidth}{!}{
\begin{tabular}{lcccc}
\toprule
变体 & NSeed & Train $R^2$ & Val $R^2$ & Val Loss \\
\midrule
A1 (MB-enrich on) & 5 & $0.9682\pm0.0087$ & $\mathbf{0.9298\pm0.0073}$ & $\mathbf{0.0670\pm0.0069}$ \\
A2 (EnrichOff-auto) & 5 & $0.9682\pm0.0087$ & $\mathbf{0.9298\pm0.0073}$ & $\mathbf{0.0670\pm0.0069}$ \\
A3 (ParasiticReplace) & 5 & $0.9694\pm0.0080$ & $0.9294\pm0.0084$ & $0.0675\pm0.0081$ \\
A4 (DualMean) & 5 & $0.9582\pm0.0216$ & $0.9234\pm0.0175$ & $0.0731\pm0.0168$ \\
A5 (MB + TopoOff) & 5 & $0.9006\pm0.0068$ & $0.8720\pm0.0034$ & $0.1222\pm0.0032$ \\
\bottomrule
\end{tabular}
}
\vspace{2pt}
\raggedright
\footnotesize
A1: T32x64 + ModeBase(enrich\_on)；
A2: T32x64 + EnrichOff(auto)；
A3: T32x64 + ParasiticReplace；
A4: T32x64 + DualMean；
A5: T32x64 + ModeBase+TopoOff。
\end{table}

主要观察如下：
\begin{enumerate}[leftmargin=1.4em]
  \item 关闭拓扑专家（TopoOff）后性能显著下降（Val $R^2$ 降低约 0.0578），说明拓扑条件残差分支是主要收益来源；
  \item \texttt{ParasiticReplace} 对主线配置未形成稳定优势，表现为中性到轻微退化；
  \item \texttt{ModeBase(enrich\_on)} 与 \texttt{EnrichOff(auto)} 在当前实现路径下数值一致，应作为等价性对照而非两条独立增益路径。
\end{enumerate}

\subsection{补充消融（TCDP(16,96)）}

为验证结论不依赖单一容量，进一步在 TCDP(16,96) 上进行同口径消融（表~\ref{tab:t16_ablation_cn}）。

\begin{table}[htbp]
\centering
\scriptsize
\setlength{\tabcolsep}{3.5pt}
\caption{TCDP(16,96) 的补充 one-factor 消融（每项 5 seeds）。}
\label{tab:t16_ablation_cn}
\resizebox{\columnwidth}{!}{
\begin{tabular}{lcccc}
\toprule
变体 & NSeed & Train $R^2$ & Val $R^2$ & Val Loss \\
\midrule
B1 (MB-enrich on) & 5 & $0.9548\pm0.0122$ & $\mathbf{0.9160\pm0.0139}$ & $\mathbf{0.0802\pm0.0132}$ \\
B2 (EnrichOff-auto) & 5 & $0.9548\pm0.0122$ & $\mathbf{0.9160\pm0.0139}$ & $\mathbf{0.0802\pm0.0132}$ \\
B3 (ParasiticReplace) & 5 & $0.9562\pm0.0118$ & $0.9151\pm0.0113$ & $0.0811\pm0.0108$ \\
B4 (DualMean) & 5 & $0.9586\pm0.0085$ & $0.9065\pm0.0168$ & $0.0892\pm0.0161$ \\
B5 (MB + TopologyOff) & 5 & $0.9006\pm0.0068$ & $0.8720\pm0.0034$ & $0.1222\pm0.0032$ \\
\bottomrule
\end{tabular}
}
\vspace{2pt}
\raggedright
\footnotesize
B1: T16x96 + mode\_base(enrich\_on)；
B2: T16x96 + enrich\_off(auto)；
B3: T16x96 + parasitic\_replace；
B4: T16x96 + dual\_readout(mean)；
B5: T16x96 + mode\_base+topology\_off。
\end{table}

该补充实验与主线结论一致：TopologyOff 退化最大；ParasiticReplace 与 DualMean 未提供稳定增益。

\end{document}

